\scanpage{69}\mspage{D}
D'ailleurs les modèles propres et normaux, pour la
\deleted{la} relation de domination, forment un
ordonné \uline{filtrant croissant}. \deleted{\unsure{On se}} Soit $U$
un \deleted{\unsure{les|tel}} modèle \inserted{régulier de $K$} \deleted{\unsure{tel l'ensemble} de ses pts
réguliers}, et soit $X_{\unsure{1|0|u}}$ \deleted{\unsure{dominé de} $X_0$} \inserted{propre \deleted{\unsure{normal}} et \uline{normal}}.

% At the upper left, three labels and an abandoned connection.
\diagram{\begin{tikzcd}[ampersand replacement=\&,row sep=small,column sep=small]
{}\&X''\arrow[dl,dash]\&{}\\
\deleted{\unsure{U}}X\arrow[rr,dash]\&{}\&X'\deleted{\unsure{X'}}
\end{tikzcd}}

\deleted{Soit $U_0$ l'ens des pts de $U_0$ où
l'application rationnelle évidente \deleted{\unsure{$X$}}
$U_0\to X_1$ induit \deleted{\unsure{un $g$}} est
définie, c'est \inserted{le compl. de \unsure{l'ens|celui} d'} \deleted{un} un fermé
codim $\geq 2$, \deleted{\unsure{donc d'où le}}
et on a donc
\deleted{complémentaire dans $X_0$ d'un ens
fermé de codim $\geq 2$.} On a un \deleted{hom
canonique}, diagramme commutatif,
d'où on voit qu'on a
un hom surjectif}

\diagram{\begin{tikzcd}[ampersand replacement=\&,column sep=small,row sep=small]
X_{\unsure{1}}\arrow[d,"f"]\&{}\&{}\\
X_0\&{}\llap{$\underset{i}{\supset}\;$}U\&{}\llap{$\underset{j}{\supset}\;$}U_{0}\arrow[ull,"g"']
\end{tikzcd}}

\diagram{\begin{tikzcd}[ampersand replacement=\&,column sep=small,row sep=small]
\pi_1(X_1)\arrow[d,"f_*"]\&{}\&\deleted{\pi_1(U_1)}\arrow[ll]\\
\pi_1(X_0)\&\pi_1(U_0)\arrow[l]\&\pi_1(U_1)\arrow[l,"\sim" ,"\substack{\text{th.}\\\text{de pureté}}"' ]\arrow[ull,"g_*"']
\end{tikzcd}}

L'application rationnelle évidente
$U\to X$ étant définie dans
le complémentaire d'une partie fermée de
codim $\geq 2$, on voit que (par le th. de
pureté) on a un hom \uline{surjectif} (car $X$ \unsure{norm.?|normal}]
\[
\pi_1(U)\to\pi_1(X)
\]
d'où à la limite un \uline{hom \uline{surjectif}}
\[
\pi_1\unsure{U;\Omega)|(U;\Omega)|(U,\Omega)}\longrightarrow
\varprojlim_{\substack{X\ \text{modèle}\\\text{propre}}}\pi_1(X;\Omega)
\]
où $\Omega$ est une extension alg. close de $K$.
\margin{\deleted{\unsure{cf}}}
% slanted
% archive sheet number: 68
